% ---------------------------------------------------------------------------
% FIGURE 1 (TikZ).  Adapted from the connectomic preprint's motif figure.
%
% Differences from that version, all forced by what this paper argues:
%   * node set is the mature circuit of this paper (I,P / L_i,L_p / INH / XOR /
%     F+,F-), not the connectomic six-node motif (I,O / EI,EO / INH / XOR);
%   * the hub returns inhibition ONLY to the two laterals, never to the
%     read-out, matching the simulated architecture in App. A.3;
%   * inhibitory edges are drawn with a filled circle head and marked (÷),
%     the near-rest shunting limit, since the text is explicit that this is an
%     idealisation of conductance-based inhibition and not its definition;
%   * the two operations the reviewer asked to separate are boxed and named:
%     the hub generates mismatch, the opponent pair supplies its sign;
%   * no Variant A/B edges: this paper commits to one wiring.
%
% Colours are those of figures/_style.py, so the vector figure and the
% matplotlib ones share one palette.
%
% Preamble requirement:
%   \usepackage{tikz}
%   \usetikzlibrary{arrows.meta,positioning,calc,backgrounds,fit}
% ---------------------------------------------------------------------------
\definecolor{excgreen}{HTML}{2E7D32}
\definecolor{inhred}{HTML}{A11A1A}
\definecolor{latblue}{HTML}{1F4E79}
\definecolor{inpgrey}{HTML}{555555}
\definecolor{outpurple}{HTML}{5A3A8A}
\definecolor{nodefill}{HTML}{F4F4F8}

\begin{figure}[t]
\centering
\begin{tikzpicture}[
  every node/.style={font=\small},
  exc/.style={-{Stealth[length=2.4mm]}, line width=0.9pt, draw=excgreen},
  drive/.style={-{Stealth[length=2.4mm]}, line width=0.9pt, draw=inpgrey},
  inh/.style={-{Circle[length=2.0mm,fill=inhred]}, line width=1.0pt, draw=inhred},
  err/.style={-{Stealth[length=2.0mm]}, line width=0.7pt, draw=black!65},
  nin/.style={circle, draw=black!55, fill=inpgrey!85, text=white,
              minimum size=7.5mm, inner sep=0pt, font=\small\bfseries},
  nlat/.style={circle, draw=black!55, fill=latblue, text=white,
               minimum size=8mm, inner sep=0pt, font=\small\bfseries},
  nhub/.style={circle, draw=black!55, fill=inhred, text=white,
               minimum size=8.5mm, inner sep=0pt, font=\scriptsize\bfseries},
  nout/.style={circle, draw=black!55, fill=outpurple, text=white,
               minimum size=8.5mm, inner sep=0pt, font=\scriptsize\bfseries},
  nfp/.style={circle, draw=black!55, fill=excgreen, text=white,
              minimum size=7.5mm, inner sep=0pt, font=\small\bfseries},
  nfm/.style={circle, draw=black!55, fill=inhred, text=white,
              minimum size=7.5mm, inner sep=0pt, font=\small\bfseries},
  zone/.style={rounded corners=3pt, draw=black!30, dashed, line width=0.6pt},
]

% ---- nodes -------------------------------------------------------------
\node[nin]  (I)   at (0,  2.05) {I};
\node[nin]  (P)   at (0, -2.05) {P};
\node[nlat] (Li)  at (2.75,  2.05) {$L_i$};
\node[nlat] (Lp)  at (2.75, -2.05) {$L_p$};
\node[nhub] (H)   at (2.75,  0)    {INH};
\node[nout] (X)   at (5.55,  0)    {XOR};
\node[nfp]  (Fp)  at (8.35,  2.05) {$F^{+}$};
\node[nfm]  (Fm)  at (8.35, -2.05) {$F^{-}$};

% ---- excitatory arms and hub recruitment -------------------------------
\draw[exc]   (I) -- (Li);
\draw[exc]   (P) -- (Lp);
\draw[drive] (I) -- (H);
\draw[drive] (P) -- (H);

% ---- the hub inhibits the laterals only, never the read-out ------------
\draw[inh] (H) -- (Li) node[midway, left=1.2mm, inhred, font=\small] {$\div$};
\draw[inh] (H) -- (Lp) node[midway, left=1.2mm, inhred, font=\small] {$\div$};

% ---- laterals pool into the read-out -----------------------------------
\draw[exc] (Li) -- (X);
\draw[exc] (Lp) -- (X);

% ---- opponent read-out: the part that supplies the sign ----------------
\draw[exc] (X)  -- (Fp);
\draw[exc] (X)  -- (Fm);
\draw[exc, bend left=16]  (Li) to (Fp);
\draw[exc, bend right=16] (Lp) to (Fm);

\node[draw=black!45, fill=nodefill, rounded corners=2pt, inner sep=3.5pt,
      font=\small] (E) at (8.35, 0) {$e=\rho(F^{+})-\rho(F^{-})$};
\draw[err] (Fp) -- (E);
\draw[err] (Fm) -- (E);

% ---- the two operations, named separately ------------------------------
\begin{scope}[on background layer]
  \node[zone, fit=(I)(P)(Li)(Lp)(H)(X), inner sep=6pt] (zA) {};
  \node[zone, fit=(Fp)(Fm)(E), inner sep=6pt] (zB) {};
\end{scope}
\node[anchor=south, font=\footnotesize\bfseries, text=black!70]
      at ($(zA.north) + (0,1.2mm)$) {coincidence suppression: \emph{mismatch}};
\node[anchor=south, font=\footnotesize\bfseries, text=black!70]
      at ($(zB.north) + (0,1.2mm)$) {opponent read-out: \emph{sign}};

% ---- operating cases ---------------------------------------------------
\node[anchor=west, align=left, font=\scriptsize, text=excgreen]
      at (10.35, 1.15) {\textbf{one input on}\\ hub weakly recruited,\\
                        lateral escapes\\ $\to$ read-out ON};
\node[anchor=west, align=left, font=\scriptsize, text=inhred]
      at (10.35, -1.15) {\textbf{both inputs on}\\ hub recruited\\
                         \emph{supralinearly}\\ $\to$ read-out OFF};

% ---- legend ------------------------------------------------------------
\node[anchor=north west, align=left, font=\scriptsize] at (-0.75, -3.35)
  {\textcolor{excgreen}{\rule[0.6mm]{4.5mm}{0.7pt}}~excitatory \quad
   \textcolor{inpgrey}{\rule[0.6mm]{4.5mm}{0.7pt}}~excitatory drive onto hub \quad
   \textcolor{inhred}{\rule[0.6mm]{4.5mm}{0.9pt}$\bullet$}~conductance-based
   inhibition ($\div$: near-rest shunting limit)};
\end{tikzpicture}

\caption{\textbf{The mature signed-XOR motif and its operations.} Each input
excites its lateral and, in parallel, the common hub, which returns
conductance-based inhibition only to the two laterals. The figure draws the
analytically clean near-rest limit as a shunt ($\div$). A single input leaves
the hub weakly recruited, so its lateral largely escapes inhibition (read-out
ON); two inputs recruit the hub \emph{supralinearly}, so that inhibition grows
faster than the drive and suppresses both laterals (read-out OFF), a soft
XOR-type gate that a downstream firing threshold makes categorical. The
read-out projects to both aggregators, and each lateral to one of them
($L_i\!\to\!F^{+}$, $L_p\!\to\!F^{-}$); the two aggregators have opposite
neurotransmitter identity, and their difference is the local signed error
$e=\rho(F^{+})-\rho(F^{-})$. The two boxed operations are distinct and are the
reason the motif is more than a gate: the central inhibition generates mismatch
selectivity, and the opponent read-out supplies its direction.}
\label{fig:mature}
\end{figure}
